\documentclass[10pt,a4paper]{amsart}
\usepackage{amsmath,amssymb,amsthm,graphicx}
% Retained preamble for the Stacks Project chapter morphisms.tex (Morphisms of Schemes).
% The chapter file pinned-original-source.tex carries only \input{preamble} on its first line, so no
% preamble text lives inside the chapter itself.  The Stacks Project uses one single shared preamble for
% every chapter; the snapshot stages it verbatim in sources/stacks-chow/preamble.tex.  The block below is
% copied mechanically, line for line, from that file's lines 14-18 (the Xy-pic setup) and lines 154-217
% (the native theorem declarations together with the mathematical macros).  Nothing was added, removed or
% redefined.  The wrapper class and the amsmath/amssymb/amsthm/graphicx base packages are supplied by the
% assembler, so the class/\externaldocument part of the shared preamble is not repeated here.
% For commutative diagrams we use Xy-pic
\usepackage[all]{xy}

% We use 2cell for 2-commutative diagrams.
\xyoption{2cell}
% Theorem environments.
%
\theoremstyle{plain}
\newtheorem{theorem}[subsection]{Theorem}
\newtheorem{proposition}[subsection]{Proposition}
\newtheorem{lemma}[subsection]{Lemma}

\theoremstyle{definition}
\newtheorem{definition}[subsection]{Definition}
\newtheorem{example}[subsection]{Example}
\newtheorem{exercise}[subsection]{Exercise}
\newtheorem{situation}[subsection]{Situation}

\theoremstyle{remark}
\newtheorem{remark}[subsection]{Remark}
\newtheorem{remarks}[subsection]{Remarks}

\numberwithin{equation}{subsection}

% Macros
%
\def\lim{\mathop{\mathrm{lim}}\nolimits}
\def\colim{\mathop{\mathrm{colim}}\nolimits}
\def\Spec{\mathop{\mathrm{Spec}}}
\def\Hom{\mathop{\mathrm{Hom}}\nolimits}
\def\Ext{\mathop{\mathrm{Ext}}\nolimits}
\def\SheafHom{\mathop{\mathcal{H}\!\mathit{om}}\nolimits}
\def\SheafExt{\mathop{\mathcal{E}\!\mathit{xt}}\nolimits}
\def\Sch{\mathit{Sch}}
\def\Mor{\mathop{\mathrm{Mor}}\nolimits}
\def\Ob{\mathop{\mathrm{Ob}}\nolimits}
\def\Sh{\mathop{\mathit{Sh}}\nolimits}
\def\NL{\mathop{N\!L}\nolimits}
\def\CH{\mathop{\mathrm{CH}}\nolimits}
\def\proetale{{pro\text{-}\acute{e}tale}}
\def\etale{{\acute{e}tale}}
\def\QCoh{\mathit{QCoh}}
\def\Ker{\mathop{\mathrm{Ker}}}
\def\Im{\mathop{\mathrm{Im}}}
\def\Coker{\mathop{\mathrm{Coker}}}
\def\Coim{\mathop{\mathrm{Coim}}}

% Boxtimes
%
\DeclareMathSymbol{\boxtimes}{\mathbin}{AMSa}{"02}

%
% Macros for moduli stacks/spaces
%
\def\QCohstack{\mathcal{QC}\!\mathit{oh}}
\def\Cohstack{\mathcal{C}\!\mathit{oh}}
\def\Spacesstack{\mathcal{S}\!\mathit{paces}}
\def\Quotfunctor{\mathrm{Quot}}
\def\Hilbfunctor{\mathrm{Hilb}}
\def\Curvesstack{\mathcal{C}\!\mathit{urves}}
\def\Polarizedstack{\mathcal{P}\!\mathit{olarized}}
\def\Complexesstack{\mathcal{C}\!\mathit{omplexes}}
% \Pic is the operator that assigns to X its picard group, usage \Pic(X)
% \Picardstack_{X/B} denotes the Picard stack of X over B
% \Picardfunctor_{X/B} denotes the Picard functor of X over B
\def\Pic{\mathop{\mathrm{Pic}}\nolimits}
\def\Picardstack{\mathcal{P}\!\mathit{ic}}
\def\Picardfunctor{\mathrm{Pic}}
\def\Deformationcategory{\mathcal{D}\!\mathit{ef}}

% --- appended 2026-10-12 (Chief): the project own definitions of the reference/slogan/history
% display environments, copied mechanically from the shared Stacks preamble lines 8-12
% (\usepackage{verbatim} plus the three \newenvironment lines, which the chapter preamble keeps
% verbatim from the project preamble). Nothing added, removed or redefined: these environments
% carry no mathematical content and the project itself typesets them as comments. Without these
% three lines a counted statement whose own text uses them fails with Environment slogan undefined.
\usepackage{verbatim}
\newenvironment{reference}{\comment}{\endcomment}
\newenvironment{slogan}{\comment}{\endcomment}
\newenvironment{history}{\comment}{\endcomment}
\usepackage{cleveref}
\crefname{originallemma}{original lemma}{original lemmas}
\Crefname{originallemma}{Original lemma}{Original lemmas}
\makeatletter
\newlabel{lemma-composition-finite-type}{{01T3}{source}{}{sourcepointer-1}{}}
\newlabel{lemma-composition-finite-type@cref}{{[originallemma][0][]01T3}{[1][0][]source}{}{}{}}
\makeatother
\crefname{originallemma}{original lemma}{original lemmas}
\Crefname{originallemma}{Original lemma}{Original lemmas}
\makeatletter
\newlabel{lemma-immersion-locally-finite-type}{{01T5}{source}{}{sourcepointer-2}{}}
\newlabel{lemma-immersion-locally-finite-type@cref}{{[originallemma][0][]01T5}{[1][0][]source}{}{}{}}
\makeatother
\crefname{originallemma}{original lemma}{original lemmas}
\Crefname{originallemma}{Original lemma}{Original lemmas}
\makeatletter
\newlabel{lemma-locally-finite-type-characterize}{{01T2}{source}{}{sourcepointer-3}{}}
\newlabel{lemma-locally-finite-type-characterize@cref}{{[originallemma][0][]01T2}{[1][0][]source}{}{}{}}
\makeatother
\crefname{originaldefinition}{original definition}{original definitions}
\Crefname{originaldefinition}{Original definition}{Original definitions}
\makeatletter
\newlabel{more-algebra-definition-J}{{07R3}{source}{}{sourcepointer-4}{}}
\newlabel{more-algebra-definition-J@cref}{{[originaldefinition][0][]07R3}{[1][0][]source}{}{}{}}
\makeatother
\title{Four equivalent formulations of the J-2 property for a locally Noetherian scheme, and the transfer to schemes locally of finite type over it}
\author{The Stacks Project Authors}
\date{stacks-project:morphisms.tex:local-snapshot-20261012; GNU FDL 1.2 or later}
\begin{document}
\maketitle
\noindent\textit{Editorial note.} All delivered mathematical text is the author's own, copied from the pinned Stacks Project chapter file pinned-original-source.tex; no statement or proof step was rewritten, reordered or removed.\par\medskip
\noindent\textit{Editorial note.} Following the multi-claim convention for this source, source\_url is the source-level URL https://stacks.math.columbia.edu/ and the source label lemma-J is carried in claim\_key and here rather than in a per-claim tag URL.\par\medskip
\noindent\textit{Editorial note.} The preamble is the chapter's shared Stacks Project preamble, retained as sources/stacks-morphisms/retained-preamble.tex: a line-for-line excerpt of the shared Stacks preamble staged with this snapshot. No macro was invented or redefined, and this card adds no preamble text of its own.\par\medskip
\noindent\textit{Editorial note.} Statement span 3309-3323 and proof span 3325-3354 come from the halfA pool card operations/pool-cards/stacks-morphisms-halfA.json and were re-read in the pinned chapter file before this card was written.\par\medskip
\noindent\textit{Editorial note.} Reread under the 2026-10-14 prefilter rule (prefilter hit is not a rejection when proof\_bytes is at least 800 or the proof runs at least 100 lines): this candidate is 1341 bytes over 30 proof lines and the prefilter matched the string 'by definition' at proof line 3337. On rereading, that occurrence is a citation of the definition of a J-2 ring (More on Algebra, Definition more-algebra-definition-J), not an author-declared omission, and the surrounding text is an argument; the candidate was therefore accepted rather than skipped.\par\medskip
\noindent\textit{Editorial note.} Tag self-check per LANE\_BRIEF 10.5 (5): the label lemma-J resolves in the project's own tag list to canonical tag 07R4 (line 9854 of the index, full label morphisms-lemma-J). The same-chapter public tag id set is empty because reference/public-items.jsonl holds no row whose primary source\_file is morphisms.tex, so the intersection of the resolved tag id with the public tag ids of this chapter is the empty set and the candidate passes. The gate was run as: python operations/agent-authoring/tools/precard\_gate.py --source morphisms.tex --label lemma-J --start 3309 --end 3323 --proof-start 3325 --proof-end 3354 --local-id 34087, giving RESULT=OK with all four span guards passing, no public interval row for this chapter, and occupancy FREE for morphisms.tex 3309 3354.\par\medskip
\noindent\textit{Editorial note.} Spans were recomputed live by label location in the pinned chapter file (locate.py) before the card was written, so the statement span 3309-3323 and the proof span 3325-3354 are the real begin/end environment lines on the current disk; the byte counts of the located text were checked against the pool candidate (statement 432 characters, proof 1135 characters after whitespace removal, matching the guard report).\par\medskip
\noindent\textit{Editorial note.} 4 cross-references occur inside the counted spans; they are listed in external\_references as navigation pointers only.\par\medskip
\noindent\textit{Allowed standard prerequisites (editorial).}\begin{itemize}
\item An open immersion is a morphism locally of finite type, cited from the same chapter. Cited at proof line 3326. Lets an open subscheme of a J-2 scheme be treated as a scheme locally of finite type over it, which is what the transfer assertion of the lemma needs.
\item A composition of morphisms locally of finite type is locally of finite type, cited from the same chapter. Cited at proof lines 3328-3329. Allows a scheme locally of finite type over an open of the base to be viewed as locally of finite type over the base itself.
\item The definition of a J-2 ring, in the form that every finite type algebra over it has open regular locus, cited from the more algebra chapter. Cited at proof lines 3335-3338. Supplies the ring-level content of the J-2 condition, which is used in both directions of the equivalence.
\item A morphism locally of finite type is characterized by finite type ring maps on affine opens, cited from the same chapter. Cited at proof lines 3347-3350. Turns local finiteness of type of the morphism into finiteness of type of the induced ring maps, so that the regular locus statement becomes available on affine pieces.
\end{itemize}\medskip
\ifdefined\hypertarget\hypertarget{sourcepointer-1}{}\fi
\noindent\textit{Source reference.} Original lemma 01T3: focused-seeds/AGP/sources/stacks-snapshot-20261011/tags.txt, source lines 2360--2360. Same-chapter reference used in the retained author text at source line(s) 2889, 2937, 3138, 3329, 9763, 9913. It is supplied only as a navigation pointer so that the retained text compiles; the counted item delivers the author's own statement and the author's complete proof, and no part of the referenced item is copied in or substituted for the counted argument.\par\medskip
\ifdefined\hypertarget\hypertarget{sourcepointer-2}{}\fi
\noindent\textit{Source reference.} Original lemma 01T5: focused-seeds/AGP/sources/stacks-snapshot-20261011/tags.txt, source lines 2362--2362. Same-chapter reference used in the retained author text at source line(s) 2980, 3135, 3326, 9944. It is supplied only as a navigation pointer so that the retained text compiles; the counted item delivers the author's own statement and the author's complete proof, and no part of the referenced item is copied in or substituted for the counted argument.\par\medskip
\ifdefined\hypertarget\hypertarget{sourcepointer-3}{}\fi
\noindent\textit{Source reference.} Original lemma 01T2: focused-seeds/AGP/sources/stacks-snapshot-20261011/tags.txt, source lines 2359--2359. Same-chapter reference used in the retained author text at source line(s) 2717, 2739, 2801, 2863, 3019, 3035, 3157, 3348, 3613, 3707, 3868, 5197, 5282, 6702, 9325, 9390, 9440, 11421, 13422, 13686. It is supplied only as a navigation pointer so that the retained text compiles; the counted item delivers the author's own statement and the author's complete proof, and no part of the referenced item is copied in or substituted for the counted argument.\par\medskip
\ifdefined\hypertarget\hypertarget{sourcepointer-4}{}\fi
\noindent\textit{Source reference.} Original definition 07R3: focused-seeds/AGP/sources/stacks-snapshot-20261011/tags.txt, source lines 9853--9853. Same-chapter reference used in the retained author text at source line(s) 3307, 3338. It is supplied only as a navigation pointer so that the retained text compiles; the counted item delivers the author's own statement and the author's complete proof, and no part of the referenced item is copied in or substituted for the counted argument.\par\medskip
\begin{lemma}
\label{lemma-J}
Let $X$ be a locally Noetherian scheme. The following are equivalent
\begin{enumerate}
\item $X$ is J-2,
\item there exists an open covering of $X$ all of whose members are
J-2 schemes,
\item for every affine open $\Spec(R) = U \subset X$ the ring
$R$ is J-2, and
\item there exists an affine open covering $S = \bigcup U_i$ such
that each $\mathcal{O}(U_i)$ is J-2 for all $i$.
\end{enumerate}
Moreover, in this case any scheme locally of finite type over $X$
is J-2 as well.
\end{lemma}

\begin{proof}
By Lemma \ref{lemma-immersion-locally-finite-type} an open immersion
is locally of finite type. A composition of morphisms locally of
finite type is locally of finite type
(Lemma \ref{lemma-composition-finite-type}). Thus it is clear that if $X$ is
J-2 then any open and any scheme locally of finite
type over $X$ is J-2 as well. This proves the final
statement of the lemma.

\medskip\noindent
If $\Spec(R)$ is J-2, then for every finite type $R$-algebra $A$
the regular locus of the scheme $\Spec(A)$ is open. Hence $R$ is
J-2, by definition (see
More on Algebra, Definition \ref{more-algebra-definition-J}).
Combined with the remarks above we conclude that (1) implies (3), and
(2) implies (4). Of course (1) $\Rightarrow$ (2) and
(3) $\Rightarrow$ (4) trivially.

\medskip\noindent
To finish the proof we show that (4) implies (1).
Assume (4) and let $Y \to X$ be a morphism locally of finite type.
We can find an affine open covering $Y = \bigcup V_j$ such that
each $V_j \to X$ maps into one of the $U_i$. By
Lemma \ref{lemma-locally-finite-type-characterize}
the induced ring map $\mathcal{O}(U_i) \to \mathcal{O}(V_j)$ is
of finite type. Hence the regular locus of
$V_j = \Spec(\mathcal{O}(V_j))$ is open. Since
$\text{Reg}(Y) \cap V_j = \text{Reg}(V_j)$ we conclude that
$\text{Reg}(Y)$ is open as desired.
\end{proof}

\end{document}
